% !TEX root = ../algebra_02.tex

\section{Изменение координат вектора при замене базиса}


\begin{Definition}{Координатный столбец}{}
	Для $v\in V$ координатный столбец $[v]_E\in K^n$ определяется равенством
	$v=E[v]_E$, то есть
	\[
	v=e_1x_1+\cdots+e_nx_n
	\quad \Longleftrightarrow \quad
	[v]_E=
	\begin{pmatrix}
		x_1\\
		\vdots\\
		x_n
	\end{pmatrix}.
	\]
\end{Definition}

Связь между координатами одного и того же вектора в разных базисах осуществляется с помощью матрицы перехода.

\begin{Theorem}{Изменение координат вектора}{}
	Пусть $E$ и $E'$ --- два базиса пространства $V$, а $C_{E'\to E}$ --- матрица перехода от $E'$ к $E$. Тогда для любого вектора $v \in V$ его координаты в базисах $E$ и $E'$ связаны соотношением:
	\[
	[v]_E = C_{E'\to E} [v]_{E'}.
	\]
\end{Theorem}

\begin{proof}
	Воспользуемся тождественным отображением $\operatorname{id}_V \colon V \to V$. Для любого вектора $v \in V$ выполнено $\operatorname{id}_V(v) = v$.
	Тогда
	\[
	[v]_E = [\operatorname{id}_V(v)]_E.
	\]
	По теореме о координатах образа вектора, матрица линейного отображения умножается на столбец координат прообраза:
	\[
	[\operatorname{id}_V(v)]_E = [\operatorname{id}_V]_{E',E} [v]_{E'}.
	\]
	Поскольку матрица перехода $C_{E'\to E}$ по определению есть $[\operatorname{id}_V]_{E',E}$, получаем искомую формулу:
	\[
	[v]_E = C_{E'\to E} [v]_{E'}.
	\]
\end{proof}
